% from ega2-1-fr.tex
\oldpage[II]{13}
<<~image directe de faisceaux quasi-cohérents~>> ne permute pas à l'opération
de <<~passage aux fibres~>>. Nous verrons cependant au chapitre~III
(\hyperref[III.4.2.4-fr]{III, 4.2.4}) un résultat dans ce sens, de nature
<<~asymptotique~>>, valable pour les faisceaux \emph{cohérents} sur $X$ lorsque
$f$ est propre (\hyperref[II.5.4-fr]{5.4}) et $S$ noethérien.
\end{remark}

\begin{corollary}[1.5.4]
\label{II.1.5.4-fr}
Pour tout préschéma $X$ affine au-dessus de $S$ et tout $s\in S$, la fibre
$f^{-1}(s)$ (où $f$ désigne le morphisme structural $X\to S$) est un schéma
affine.
\end{corollary}

Il suffit d'appliquer (\hyperref[II.1.5.1-fr]{1.5.1}) à
$S'=\operatorname{Spec}(k(s))$ et d'utiliser
(\hyperref[II.1.3.4-fr]{1.3.4}).

\begin{corollary}[1.5.5]
\label{II.1.5.5-fr}
Soient $X$ un $S$-préschéma, $S'$ un préschéma affine au-dessus de $S$~; alors
$X'=X_{(S')}$ est un préschéma affine au-dessus de $X$. En outre, si
$f:X\to S$ est le morphisme structural, il y a un isomorphisme canonique de
$\mathcal{O}_X$-Algèbres
$\mathcal{A}(X')\xrightarrow{\sim}f^*(\mathcal{A}(S'))$, et pour tout
$\mathcal{A}(S')$-Module quasi-cohérent $\mathcal{M}$, un di-isomorphisme
canonique
$f^*(\mathcal{M})\xrightarrow{\sim}
\mathcal{A}(f^{\prime*}(\widetilde{\mathcal{M}}))$, en désignant par
$f'=f_{(S')}$ le morphisme structural $X'\to S'$.
\end{corollary}

Il suffit d'intervertir les rôles de $X$ et de $S'$ dans
(\hyperref[II.1.5.1-fr]{1.5.1}) et
(\hyperref[II.1.5.2-fr]{1.5.2}).

\begin{env}[1.5.6]
\label{II.1.5.6-fr}
Soient maintenant $S$, $S'$ deux préschémas, $q:S'\to S$ un morphisme,
$\mathcal{B}$ (resp. $\mathcal{B}'$) une $\mathcal{O}_S$-Algèbre (resp. une
$\mathcal{O}_{S'}$-Algèbre) quasi-cohérente, $u:\mathcal{B}\to\mathcal{B}'$
un $q$-morphisme (c'est-à-dire un homomorphisme
$\mathcal{B}\to q_*(\mathcal{B}')$ de $\mathcal{O}_S$-Algèbres). Si
$X=\operatorname{Spec}(\mathcal{B})$,
$X'=\operatorname{Spec}(\mathcal{B}')$, on en déduit canoniquement un
morphisme
\[
  v=\operatorname{Spec}(u):X'\to X
\]
tel que le diagramme
\[
  \xymatrix{
    X'\ar[r]^v\ar[d] & X\ar[d]\\
    S'\ar[r]_q & S
  }
  \tag{1.5.6.1}\label{II.1.5.6.1-fr}
\]
soit commutatif (les flèches verticales étant les morphismes structuraux). En
effet, la donnée de $u$ équivaut à celle d'un homomorphisme de
$\mathcal{O}_{S'}$-Algèbres quasi-cohérentes
$u^\sharp:q^*(\mathcal{B})\to\mathcal{B}'$
(\hyperref[0.4.4.3-fr]{0, 4.4.3})~; elle définit donc canoniquement un
$S'$-morphisme
\[
  w:\operatorname{Spec}(\mathcal{B}')\to
  \operatorname{Spec}(q^*(\mathcal{B}))
\]
tel que $\mathcal{A}(w)=u^\sharp$
(\hyperref[II.1.2.7-fr]{1.2.7}). D'autre part, il résulte de
(\hyperref[II.1.5.2-fr]{1.5.2}) que
$\operatorname{Spec}(q^*(\mathcal{B}))$ s'identifie canoniquement à
$X\times_S S'$~; le morphisme $v$ est le composé
$X'\xrightarrow{w}X\times_S S'\xrightarrow{p_1}X$ de $w$ avec la première
projection, et la commutativité de
(\hyperref[II.1.5.6.1-fr]{1.5.6.1}) résulte des définitions. Soient $U$ (resp.
$U'$) un ouvert affine de $S$ (resp. $S'$) tels que $q(U')\subset U$,
$A=\Gamma(U,\mathcal{O}_S)$, $A'=\Gamma(U',\mathcal{O}_{S'})$ leurs anneaux,
$B=\Gamma(U,\mathcal{B})$, $B'=\Gamma(U',\mathcal{B}')$~; la restriction de
$u$ en un $(q|U')$-morphisme
$\mathcal{B}|U\to\mathcal{B}'|U'$ correspond à un di-homomorphisme d'algèbres
$B\to B'$~; si $V$, $V'$ sont les images réciproques de $U$, $U'$ dans $X$,
$X'$ respectivement, par les morphismes structuraux, le morphisme $V'\to V$,
restriction de $v$, correspond (\hyperref[I.1.7.3-fr]{I, 1.7.3}) au
di-homomorphisme précédent.
\end{env}

\begin{env}[1.5.7]
\label{II.1.5.7-fr}
Sous les mêmes hypothèses que dans (\hyperref[II.1.5.6-fr]{1.5.6}), soit
$\mathcal{M}$ un $\mathcal{B}$-Module quasi-cohérent~; il y a alors un
isomorphisme canonique de $\mathcal{O}_{X'}$-Modules
\[
  v^*(\widetilde{\mathcal{M}})\xrightarrow{\sim}
  \bigl(q^*(\mathcal{M})\otimes_{q^*(\mathcal{B})}\mathcal{B}'\bigr)^{\sim}
  \tag{1.5.7.1}\label{II.1.5.7.1-fr}
\]

